% ---------------------------------------------------------------------
%  CHƯƠNG 4 — CÀI ĐẶT VÀ THỰC NGHIỆM
% ---------------------------------------------------------------------
\chapter{Cài đặt và thực nghiệm}
\label{ch:caidat}

\begin{muctieu}
Chương này mô tả môi trường công nghệ triển khai, cách thức hiện thực hóa
các dịch vụ cốt lõi trong hệ thống Hybrid AI, kiến trúc xử lý công việc bất
đồng bộ và minh hoạ các giao diện người dùng thực tế đã được chụp từ hệ thống.
\end{muctieu}

\section{Môi trường công nghệ triển khai}

Hệ thống được phát triển và vận hành thực nghiệm trên máy trạm với cấu hình:
\begin{itemize}[leftmargin=1.6em, itemsep=2pt]
  \item \textbf{Phần cứng}: CPU AMD Ryzen 7 / Intel Core i7, 16~GB RAM, ổ cứng SSD NVMe.
  \item \textbf{Hệ điều hành}: Microsoft Windows 11 Pro 64-bit, môi trường thực
        thi dòng lệnh Git Bash (POSIX shell).
  \item \textbf{Ngôn ngữ và Framework máy chủ}: PHP 8.4 (yêu cầu tối thiểu 8.2),
        Laravel Framework 13.9 sử dụng kiến trúc MVC và Eloquent ORM.
  \item \textbf{Cơ sở dữ liệu}: MySQL 8.0 với bảng mã UTF-8 đầy đủ (\code{utf8mb4}),
        lập chỉ mục kết hợp trên các trường trạng thái và độ ưu tiên.
  \item \textbf{Giao diện người dùng}: Blade Template Engine, Tailwind CSS 3.1,
        Vite 6.0 và Chart.js 4.4.
  \item \textbf{Cổng giao tiếp AI}: Dịch vụ tương thích OpenAI Chat Completions
        API với mô hình đa phương thức, hỗ trợ gọi công cụ (\code{function
        calling}) và cơ chế chuyển đổi dự phòng giữa nhà cung cấp chính/phụ.
\end{itemize}

\section{Hiện thực hóa các dịch vụ cốt lõi}

Toàn bộ logic nghiệp vụ được đóng gói trong các lớp Service chuyên biệt, đặt
tại thư mục \code{app/Services/}.

\subsection{Dịch vụ trích xuất sự kiện \code{TriageAiFactService}}

Dịch vụ này gửi chuỗi mô tả của người bệnh tới mô hình ngôn ngữ cùng lược đồ
công cụ \code{extract\_triage\_facts}. Danh sách mã triệu chứng được giới hạn
chặt chẽ bằng tham số \code{\$allowedCodes} lấy từ bảng
\code{medical\_symptoms}. Trước khi gọi mô hình, phương thức \code{knownFacts}
tự nạp các dữ kiện hồ sơ người dùng đã khai (tuổi, thai kỳ, cho con bú, đã
dùng thuốc) để đưa vào lời nhắc, tránh hỏi lại người bệnh những thông tin hệ
thống đã có. Toàn bộ gọi mô hình nằm trong khối \code{try/catch}: mọi lỗi
nhà cung cấp chỉ được ghi log và trả về \code{null} để luồng luật tất định
tiếp tục, như Mã nguồn~\ref{lst:fact} minh hoạ.

\begin{lstlisting}[language=PHP, caption={Trích xuất sự kiện lâm sàng qua Tool-calling.}, label={lst:fact}]
public function extract(
    TriageSession $session,
    string $text,
    array $allowedCodes,
    array $wizardTranscript = []
): ?array {
    if (! $this->isEnabled() || trim($text) === '') {
        return null;
    }

    try {
        $known = $this->knownFacts($session, $text);

        $response = $this->callProvider($this->provider(), [
            ['role' => 'system', 'content' => $this->systemPrompt()],
            ['role' => 'user', 'content' => $this->buildUserMessage($session, $text, $wizardTranscript, $known['prompt'])],
        ], [
            'functions' => [$this->toolSchema($allowedCodes)],
            'function_call' => 'auto',
        ]);

        $raw = $this->decodeToolCall($response);

        if ($raw === null) {
            return null;
        }

        return $this->validate($raw, $allowedCodes, $known['topics'],
            $this->deniedSymptomTerms($session));
    } catch (\Throwable $exception) {
        Log::warning('Triage AI fact extraction failed');

        return null;
    }
}
\end{lstlisting}

\subsection{Động cơ suy diễn bảo thủ \code{MedicalTriageService}}

Trái tim của hệ thống là phương thức \code{resolveOutcomeWithAi}, trình bày
tại Mã nguồn~\ref{lst:evaluate}. Động cơ thực hiện đánh giá hai lần: một lần
trên các sự kiện rút trích từ biểu thức chính quy cơ bản (\code{\$baseRule})
và một lần trên dữ liệu hợp nhất với AI (\code{\$aiRule}).

\begin{lstlisting}[language=PHP, caption={Cơ chế phòng vệ kép đảm bảo an toàn lâm sàng.}, label={lst:evaluate}]
protected function resolveOutcomeWithAi(
    TriageSession $session,
    string $text,
    Collection $symptoms,
    array $facts
): array {
    $baseRule = $this->resolveRule($symptoms, $facts);
    $factService = app(TriageAiFactService::class);

    if (!$factService->isEnabled()) {
        return ['rule' => $baseRule, 'symptoms' => $symptoms, 'ai' => null];
    }

    $aiFacts = $factService->extract($session, $text,
        $this->symptomDictionary()->pluck('code')->all(),
        $this->wizardTranscript($session)
    );

    if ($aiFacts === null) {
        return ['rule' => $baseRule, 'symptoms' => $symptoms, 'ai' => null];
    }

    $aiSymptoms = $symptoms
        ->concat($this->symptomDictionary()->whereIn('code', $aiFacts['symptom_codes']))
        ->unique('MaTC')->values();

    $aiRule = $this->resolveRule($aiSymptoms,
        $factService->mergeConservatively($facts, $aiFacts));

    // |Nguyên tắc then chốt: Giữ lại luật có mức độ nghiêm trọng hơn|
    $winner = $this->severerRule($baseRule, $aiRule);

    return [
        'rule' => $winner,
        'symptoms' => $aiSymptoms,
        'ai' => ['status' => 'ok', 'winning_outcome' => $winner?->outcome]
    ];
}
\end{lstlisting}

\subsection{Dịch vụ diễn giải và bộ kiểm duyệt \code{TriageAiBriefingService}}

Dịch vụ \code{TriageAiBriefingService} chịu trách nhiệm soạn thảo văn bản giải
thích chi tiết sau khi kết luận phân tầng đã được xác định. Điểm cốt lõi là
bộ kiểm duyệt \code{sanitize} tại tầng ứng dụng (Mã nguồn~\ref{lst:sanitize}):

\begin{lstlisting}[language=PHP, caption={Bộ kiểm duyệt loại bỏ các kết quả AI không đạt chuẩn.}, label={lst:sanitize}]
protected function sanitize(array $raw, array $candidates, TriageSession $session): ?array
{
    $why = $this->cleanWhy($raw['why'] ?? '');

    $briefing = ['why' => $why];

    foreach (self::SECTIONS as $section) {
        $briefing[$section] = $this->cleanList($raw[$section] ?? null);
    }

    // |Id sản phẩm được giao với danh mục ứng viên từ catalogue,
    // mô hình không thể gán tên hay giá thuốc tùy ý|
    $briefing['products'] = $this->cleanProducts($raw['products'] ?? null,
        $candidates, $session);

    // |Loại bỏ hoàn toàn nếu thiếu lý do hoặc không đủ 2/5 danh mục hướng dẫn|
    $filledSectionCount = collect(self::SECTIONS)
        ->filter(fn (string $s) => $briefing[$s] !== [])
        ->count();
    $hasContent = $why !== '' && $filledSectionCount >= 2;

    return $hasContent ? $briefing : null;
}
\end{lstlisting}

\section{Kiến trúc bất đồng bộ và cơ chế tự phục hồi}

Để không làm tắc nghẽn yêu cầu HTTP của người dùng, hệ thống sử dụng kiến trúc
xử lý nền gồm ba chặng:

\begin{enumerate}[leftmargin=1.6em, itemsep=2pt]
  \item \textbf{Đẩy tác vụ vào hàng đợi}: Khi người dùng hoàn thành bước cuối
        của quy trình sàng lọc, hệ thống ghi bản ghi phiên với trạng thái
        \code{completed}, đồng thời tạo một bản ghi \code{AiJob} thuộc loại
        \code{triage\_briefing} được xếp vào hàng đợi nền
        (\code{ProcessTriageBriefing}).
  \item \textbf{Thăm dò giao diện (Client-side Polling)}: Trình duyệt nhận ngay
        trang kết quả hiển thị nhóm nguy cơ và lời khuyên chuẩn. Khối giao diện
        diễn giải hiển thị hiệu ứng khung xương tải trước (\emph{skeleton loader})
        và gửi yêu cầu kiểm tra trạng thái \code{AiJob} theo lịch giãn dần
        (900--2~400~ms, tối đa 60 lượt) thay vì một chu kỳ cố định.
  \item \textbf{Cơ chế cứu hộ đồng bộ (Synchronous Fallback Seam)}: Một tác vụ
        vẫn ở trạng thái \code{queued} sau 12 giây mà không có worker nào đang
        chạy (\code{AiJob::isStranded} kiểm tra qua \code{hasLiveWorker}) được
        coi là bị bỏ rơi. Khi người dùng tải lại trang, phương thức
        \code{ensureBriefingJob} trong \code{HealthCheckController} phát hiện và
        thực thi đồng bộ ngay trong lượt tải lại đó, bảo đảm người dùng không
        phải chờ vô hạn. Ca cấp cứu không bao giờ vào hàng đợi: cờ đỏ được phát
        hiện bằng PHP thuần nên kết luận trả về ngay trong yêu cầu.
\end{enumerate}

\section{Giao diện thực tế của hệ thống}

Các hình ảnh dưới đây được chụp trực tiếp từ ứng dụng thực tế thông qua kịch
bản tự động hóa kiểm thử Playwright trên cổng thử nghiệm cục bộ.

\subsection{Giao diện khảo sát triệu chứng từng bước}

Hình~\ref{fig:wizard} minh hoạ bước đầu tiên trong quy trình khảo sát 5 bước.
Người dùng có thể chọn nhanh các triệu chứng thường gặp bằng thẻ nút bấm hoặc
nhập mô tả tự do bằng tiếng Việt, kèm cảnh báo pháp lý về tính chất sàng lọc y tế.

\begin{figure}[H]
\centering
\includegraphics[width=0.92\textwidth]{01-wizard-step1.png}
\caption{Giao diện khảo sát triệu chứng trực quan với thanh tiến trình và các nút chọn nhanh}
\label{fig:wizard}
\end{figure}

\subsection{Giao diện kết quả phân tầng tự chăm sóc kèm AI Briefing}

Hình~\ref{fig:result_briefing} thể hiện trang kết quả khi hệ thống xếp ca bệnh
vào nhóm \term{Có thể tự chăm sóc}. Bên cạnh kết luận chính, khối
\emph{«Em giải thích kỹ hơn cho anh/chị»} được gắn nhãn AI minh bạch, hiển thị
lý do lâm sàng dựa trên thông tin người bệnh đã kể và các thẻ gợi ý thuốc không
kê đơn phù hợp.

\begin{figure}[H]
\centering
\includegraphics[width=0.82\textwidth]{03-result-briefing.png}
\caption{Trang kết quả phân tầng tự chăm sóc hiển thị phần diễn giải chi tiết AI Briefing}
\label{fig:result_briefing}
\end{figure}

\subsection{Giao diện kết quả cảnh báo cấp cứu khẩn cấp}

Hình~\ref{fig:result_emergency} thể hiện phản ứng tức thời của hệ thống khi
phát hiện cờ đỏ cấp cứu (ví dụ: đau ngực dữ dội, khó thở). Hệ thống lập tức
bật viền đỏ cảnh báo, loại bỏ toàn bộ phần gợi ý thuốc và hiển thị nổi bật nút
bấm gọi cấp cứu 115, không cho phép tiếp tục tự chăm sóc.

\begin{figure}[H]
\centering
\includegraphics[width=0.88\textwidth]{04-result-emergency.png}
\caption{Giao diện cảnh báo cấp cứu lập tức kích hoạt khi phát hiện cờ đỏ nguy kịch}
\label{fig:result_emergency}
\end{figure}

\subsection{Giao diện hàng đợi duyệt ca của dược sĩ}

Hình~\ref{fig:pharmacist_queue} là phân hệ nội bộ dành cho dược sĩ chuyên môn.
Các ca bệnh thuộc nhóm \code{pharmacist\_review} hoặc do người bệnh bấm yêu
cầu hỗ trợ sẽ được đưa vào hàng đợi tập trung, cho phép dược sĩ tiếp nhận, xem
lịch sử triệu chứng và gửi phản hồi y khoa chính thức.

\begin{figure}[H]
\centering
\includegraphics[width=0.92\textwidth]{05-pharmacist-queue.png}
\caption{Giao diện hàng đợi phân tầng dành cho dược sĩ tiếp nhận và thẩm định ca bệnh}
\label{fig:pharmacist_queue}
\end{figure}
